\documentclass[11pt]{article}

\usepackage[margin=1in]{geometry}
\usepackage{amsmath,amssymb,amsthm,mathtools}
\usepackage{enumitem}
\usepackage{hyperref}
\hypersetup{colorlinks=true,linkcolor=blue,citecolor=blue,urlcolor=blue}
\usepackage{orcidlink}
\usepackage{fontawesome5}

\newtheorem{theorem}{Theorem}[section]
\newtheorem{proposition}[theorem]{Proposition}
\newtheorem{lemma}[theorem]{Lemma}
\newtheorem{corollary}[theorem]{Corollary}
\newtheorem{empirical}[theorem]{Empirical Result}
\theoremstyle{remark}
\newtheorem{remark}[theorem]{Remark}

\DeclareMathOperator{\ord}{ord}

\title{A Calibrated Separation of the Kontorovich--Lagarias\\
and Volkov Exponent Predictions for the $5x+1$ Reverse Tree}

\author{Renato Augusto Tavares~{\small\href{mailto:dr.renatotavares@gmail.com}{\textcolor{black}{\faEnvelope}}}~\orcidlink{0009-0002-0196-3311}\\
\normalsize Universidade Federal de Goiás\\
\small \href{https://orcid.org/0009-0002-0196-3311}{https://orcid.org/0009-0002-0196-3311}}
\date{\today}

\begin{document}

\maketitle

\begin{abstract}
% DRAFT -- rewrite by hand before submission (Rule 4b: abstracts carry
% the highest marker density of any section in the paper).
Kontorovich and Lagarias predict a counting exponent
$\eta_{5,BP}\approx0.650919$ for the reverse tree of the accelerated
$5x+1$ map, and note a competing prediction $\eta^\ast_{5,BP}\approx
0.678$ due to Volkov, leaving the two unresolved for lack of
sufficiently powerful data. We report an exact enumeration with
Aitken-extrapolated truncation correction that reaches only $0.639$,
a fixed-window reading that excludes neither value. We then show the
estimator itself is biased by more than the gap separating the two
predictions when applied to a process of known exponent, so a raw
reading cannot settle the dispute at all. Building synthetic controls
that share the arithmetic tree's branching and sibling structure but
target a chosen exponent, via an exact sibling congruence in the
reverse tree, we obtain a calibrated comparison at the depth where
the controls read their own target exponent to within $0.003$: the
arithmetic tree separates clearly from a control at $0.678$ and
matches controls at $0.650919$. The result is a measurement with
calibrated controls, not a proof, and targets the exponent value in
dispute rather than Volkov's branching model itself.
\end{abstract}

\section{Introduction}
\label{sec:intro}

The Collatz, or $3x+1$, conjecture asks whether every positive integer
eventually reaches $1$ under the map $n\mapsto n/2$ ($n$ even),
$n\mapsto3n+1$ ($n$ odd). Kontorovich and Lagarias
\cite{KontorovichLagarias2009} study the generalization to $qx+1$ maps
through the counting statistics of the associated reverse tree, and
give a stochastic model predicting a counting exponent
$\eta_{q,BP}$ governing the growth of $N_u(x)$, the number of integers
up to $x$ that eventually map to a fixed residue $u$. At $q=5$ this
model predicts $\eta_{5,BP}\approx0.650919$. In the same reference the
authors discuss a competing branching model due to Volkov, predicting
$\eta^\ast_{5,BP}\approx0.678$, and note explicitly that the two
predictions are close enough ($\Delta=0.027$) that available numerical
data could not discriminate between them, calling for further
investigation.

This paper reports such an investigation. \S\ref{sec:setup} fixes the
reverse-tree construction and records an exact sibling congruence used
later to build calibrated controls. \S\ref{sec:dispute} restates the
dispute precisely. \S\ref{sec:method}--\S\ref{sec:result} report a
direct exact enumeration, reaching a fixed-window reading of $0.639$
that, on its own, excludes neither prediction.
\S\ref{sec:calibrated} is the paper's contribution: we show this
estimator is biased on a process of \emph{known} exponent by more than
the gap in dispute, so the raw reading in \S\ref{sec:result} was never
evidence against either side, and we then build matched synthetic
controls that isolate and calibrate out that bias, obtaining a clear
separation from the disputed value $0.678$.

\section{Setup: the reverse tree of the accelerated $qx+1$ map}
\label{sec:setup}

Fix an odd integer $q\ge3$ coprime to $2$. The \emph{accelerated
$qx+1$ map} is
\[
  T_q(n) = \frac{qn+1}{2^{v_2(qn+1)}}, \qquad n \text{ odd},
\]
where $v_2$ is the $2$-adic valuation. The \emph{reverse tree} rooted
at an odd integer $u$ coprime to $q$ is built by, at each odd node
$u$, adjoining a child $w_a := (2^a u - 1)/q$ for every $a\ge1$ such
that $2^a u\equiv1\pmod q$; such $w_a$ is automatically an odd integer
whenever it is an integer. Writing $d:=\ord_q(2)$, admissibility of
$a$ for a given parent $u$ depends only on $u\bmod q$: if some
$a_0(u)\in\{1,\dots,d\}$ satisfies $2^{a_0(u)}u\equiv1\pmod q$, it is
unique, and the admissible exponents are exactly $a\equiv a_0(u)\pmod
d$. At $q=5$, $d=4$ and $A_0=\{1\mapsto4,\allowbreak2\mapsto3,
\allowbreak3\mapsto1,\allowbreak4\mapsto2\}$; nodes with $u\equiv
0\pmod5$ are sterile, and, unlike $q=3$, there is no further parity
condition beyond integrality mod $5$.

We count distinct integer labels rather than path occurrences:
\[
 N_u(x):=\#\{n\le x:n\text{ odd and }T_q^k(n)=u
                    \text{ for some }k\ge0\}.
\]

\begin{lemma}[Sibling congruence]
\label{lem:sibling-congruence}
Let $q\ge3$ be odd, $d=\ord_q(2)$, and $u$ odd with $2^au\equiv1\pmod
q$ solvable. The children of $u$ in the accelerated reverse tree are
$w_j=(2^{a_0+jd}u-1)/q$, $j\ge0$, where $a_0=a_0(u)$. Then
\[
  w_{j+1} = 2^d w_j + \frac{2^d-1}{q} \qquad\text{(exact identity in
  $\mathbb Z$)},
\]
and consequently $w_{j+1}\equiv w_j+c\pmod q$, with
$c:=\bigl((2^d-1)/q\bigr)\bmod q$.
\end{lemma}

\begin{proof}
Direct substitution: $2^d w_j+(2^d-1)/q =
\bigl(2^{a_0+(j+1)d}u-2^d+2^d-1\bigr)/q = w_{j+1}$. The congruence
follows from $2^d\equiv1\pmod q$.
\end{proof}

When $q$ is prime, $c=0$ would force $q^2\mid2^d-1$, the base-$2$
Wieferich condition; the only known Wieferich primes are $1093$ and
$3511$, so $c\ne0$ for every $q$ of interest here, and consecutive
siblings' residues mod $q$ advance by a fixed nonzero step. At $q=5$,
$d=4$, $c=3$.

\section{The dispute}
\label{sec:dispute}

For the reverse tree of $T_5$, Kontorovich--Lagarias
\cite{KontorovichLagarias2009} predict a counting exponent
$\eta_{5,BP}\approx0.650919$, while a competing stochastic model due
to Volkov, discussed in the same reference, predicts
$\eta^\ast_{5,BP}\approx0.678$. The authors themselves flag this as an
open dispute, unresolved for lack of sufficiently powerful data; the
two predictions differ by only $\Delta=0.027$.

\section{Method}
\label{sec:method}

We implemented the exact reverse tree of $T_5$ from scratch (the
admissibility rule of \S\ref{sec:setup} specialized to $q=5$).

A first pilot ($60$ roots) already gave a slope of $0.6433$ (95\% CI
$[0.6327,0.6531]$), excluding $0.678$ by roughly $6.4$ standard
errors. A larger production run ($n=300$ roots) initially sampled
roots too close to the measurement window's lower checkpoint,
contaminating the truncation buffer; correcting this (roots restricted
well below the window) gave a slope of $0.63801$ (95\% CI
$[0.63159,0.64410]$) --- which, at this sample size, excluded
\emph{both} candidate values ($0.678$ by $12.6$ standard errors, but
also $0.650919$ by $4.1$ standard errors), a failure mode anticipated
in advance: once the standard error falls below the residual
truncation bias, a naive confidence interval on the biased estimator
excludes everything, which is not evidence against the correct
prediction.

We therefore rewrote the enumeration to track, in a single pass, the
running path-maximum at every node, giving simultaneous counts at
truncation buffers $10^9$ through $10^{13}$ (validated byte-for-byte
against the earlier buffer-by-buffer method on four test roots). The
pooled mean curve across roots decays geometrically:
\[
  0.60049\ (10^9) \to 0.62387\ (10^{10}) \to 0.63261\ (10^{11}) \to
  0.63650\ (10^{12}) \to 0.63801\ (10^{13}),
\]
with increments $0.0234, 0.0087, 0.0039, 0.0015$ shrinking by a factor
of roughly $0.4$--$0.5$ per decade of truncation --- a signature of a
genuine, extrapolable truncation bias, not noise.

\section{Result: a fixed-window estimator}
\label{sec:result}

Aitken's $\Delta^2$ extrapolation of this sequence to infinite
truncation gives
\begin{equation}\label{eq:kl-result}
  \hat\eta_5 = 0.639, \qquad 95\%\ \text{bootstrap interval, fixed window} =
  [0.633,\, 0.645].
\end{equation}
The bootstrap interval measures variation between sampled roots after
the truncation extrapolation. It does not include the remaining bias
from the fixed measurement window and therefore cannot exclude either
asymptotic prediction. That bias has a visible signature: the
slope-per-decade panel \emph{within} the fixed measurement window
($10^4\!\to\!10^5$ through $10^7\!\to\!10^8$) is still monotonically
rising at the last decade tested ($0.6021\to0.6296\to0.6432\to0.6460$,
no plateau), approaching $0.6509$ monotonically without having reached
it --- the signature of \emph{fixed-window pre-asymptotics} (the
window is not yet ``deep enough''), not of uncorrected truncation bias
(already handled by the Aitken $\Delta^2$ extrapolation above).

\begin{empirical}[Finite-window comparison of two exponent predictions]
\label{thm:kl}
The fixed-window estimator of the $5x+1$ reverse-tree counting
exponent, computed by exact enumeration with
Aitken-$\Delta^2$-extrapolated truncation correction, is $0.639$
(bootstrap interval $[0.633,0.645]$). The sequence of local slopes
rises monotonically toward the Kontorovich--Lagarias prediction
$0.650919$ on the available windows. Since it has not stabilized, the
experiment does not provide a calibrated confidence interval for the
asymptotic exponent and does not exclude the Volkov prediction.
\end{empirical}

\section{A calibrated comparison}
\label{sec:calibrated}

The estimator behind Empirical Result~\ref{thm:kl} had never been run
on a process with a known exponent. It is biased: on a branching
random walk with the same offspring law but an independently drawn
branch class at every node (exponent provably $0.650919$, the
Kontorovich--Lagarias prediction itself, realized as a pressure-equation
root in a companion paper on the $qx+1$ branching structure
\cite{PressureCompanion}), the same window, truncation, and
extrapolation reads $0.6131$, an under-read of $0.038$, larger than
the gap $\Delta=0.027$ separating the two disputed values. Comparing
the raw reading against either theoretical prediction therefore
settles nothing.

The bias depends on how much a process fluctuates, and it can be
calibrated out by comparison rather than subtraction: run the same
estimator on synthetic processes built to have exponent $0.650919$ and
exponent $0.678$ respectively, sharing every feature of the arithmetic
tree except the ultimate value of $q$, and see which reading the
arithmetic tree matches. Lemma~\ref{lem:sibling-congruence} makes this
construction exact: replacing the integer denominator $q$ by a tunable
real value reproduces the arithmetic tree's branching and its sibling
spacing while moving the exponent to any target solving
$q_{\mathrm{val}}^{\alpha}=q(2^{\alpha}-1)$.

\begin{empirical}[Calibrated separation from the Volkov exponent]
\label{thm:kl-calibrated}
At truncation depth $10^9\!\to\!10^{10}$, where the calibrated controls
read their target exponent to within $0.003$, the arithmetic tree
reads $0.64926$ ($95\%$ interval $[0.64818,0.65027]$); a control built
to exponent $0.650919$ reads $0.64981$ and $0.65122$ on two
independent constructions; a control built to exponent $0.678$ reads
$0.67748$, ten interval-widths from the arithmetic reading. The
separation widens, not narrows, as bias is calibrated away with
depth. This measurement compares against the exponent value $0.678$,
not against Volkov's branching model itself, which has a different
tree structure and is not implemented here.
\end{empirical}

\section{Discussion}
\label{sec:discussion}

Two systematics were checked rather than assumed before trusting
Empirical Result~\ref{thm:kl-calibrated}. The integer-denominator
construction and its real-valued counterpart at the exponent
$0.650919$ agree to $0.0035$ over six seeds ($1.5$ standard errors),
bounding any implementation systematic between the two recursions at
about $0.004$. And the exponent-$0.678$ control fluctuates more than
the arithmetic tree, which would if anything bias its reading
\emph{downward} toward the arithmetic reading; the separation survives
that direction of bias.

This is a measurement with calibrated controls, not a proof. It settles
the narrower, well-defined question of whether the arithmetic tree's
counting exponent is closer to $0.650919$ or to the value $0.678$;
it does not implement or test Volkov's branching model directly, which
uses a different encoding of the iterates than the matched controls
built here. Whether that stronger claim holds is a separate question.

\section*{Data Availability Statement}

Code and data reproducing every claim in this paper, organized by
section with a README per subfolder, are at
\url{https://github.com/faculdade/collatz-kl-volkov}.

\section*{Acknowledgments}

AI assistance (Claude, Anthropic) was used for textual review and
translation.

\bibliographystyle{plain}
\begin{thebibliography}{9}

\bibitem{KontorovichLagarias2009} A.\ V.\ Kontorovich and J.\ C.\
  Lagarias, \emph{Stochastic models for the $3x+1$ and $5x+1$
  problems}, in: \emph{The Ultimate Challenge: the $3x+1$ Problem}
  (J.\ C.\ Lagarias, ed.), American Mathematical Society, 2010.

\bibitem{PressureCompanion} R.\ A.\ Tavares, \emph{A closed-form
  pressure equation for the accelerated $qx+1$ branching process},
  companion paper, in preparation.

\end{thebibliography}

\end{document}
